\chapter{Troubleshooting}
\label{sec:man-troubleshooting}

\section{401 Unauthorized}

\textbf{Síntoma}: la UI o el MCP reciben 401.

\textbf{Causa}: token ausente o desalineado.

\textbf{Fix}:

\begin{enumerate}
    \item En los logs del app buscar la línea
    \code{GRAPHHUNTER_API_TOKEN=<uuid>}.
    \item Copiar el UUID al cliente (archivo \filepath{.env},
    config MCP, CLI \code{--token}).
    \item Reiniciar el cliente.
\end{enumerate}

\section{Connection refused}

\textbf{Síntoma}: el cliente no logra conectarse.

\textbf{Causa}: app no corre, puerto distinto, o interfaz de bind
incorrecta.

\textbf{Fix}:

\begin{itemize}
    \item Verificar que el app esté abierto. La API corre sólo
    mientras el app está vivo.
    \item En Windows, preferir \code{127.0.0.1} sobre
    \code{localhost} en \code{GRAPHHUNTER_API_URL}.
    \item Si cambiaste el puerto con
    \code{GRAPHHUNTER_API_PORT}, actualizar el cliente.
    \item Si configuraste \code{GRAPHHUNTER_API_HOST=0.0.0.0},
    revisar firewall.
\end{itemize}

\section{HTML en lugar de JSON}

\textbf{Síntoma}: el cliente recibe HTML donde esperaba JSON.

\textbf{Causa}: otro servicio ocupa el puerto 37891 (común si
tenés un dev server por ahí).

\textbf{Fix}: parar el servicio conflictivo, o cambiar el puerto
del app con \code{GRAPHHUNTER_API_PORT=NN} y sincronizar el
cliente.

\section{Out of memory con archivos grandes}

\textbf{Síntoma}: ingesta falla o el app se congela con EVTX de
varios GB.

\textbf{Fix}:

\begin{itemize}
    \item Activar \emph{Ingesta por disco}
    (cap.~\ref{sec:man-ingesta}).
    \item Reducir \code{GRAPHHUNTER_INGEST_BATCH_SIZE} (10000 o
    5000).
    \item Partir el archivo antes e ingestar por partes.
\end{itemize}

\section{MCP client no ve tools nuevas}

\textbf{Síntoma}: agregaste una tool en el MCP pero el cliente no
la lista.

\textbf{Fix}: bajá y subí el cliente MCP (Cursor, Claude Code,
etc.). Los clientes cachean la lista de tools al arrancar.

\section{Datos demo fallan}

\textbf{Síntoma}: al ingestar un archivo de \filepath{demo_data/},
error de path o de formato.

\textbf{Fix}:

\begin{itemize}
    \item Verificar que arrancaste el app desde el directorio
    raíz del proyecto, no desde \filepath{apps/tauri/}.
    \item Si descargaste los datos reales
    (\filepath{DOWNLOAD_REAL_DATA.md}), verificar que estén bajo
    \filepath{demo_data/} con las subcarpetas esperadas.
\end{itemize}

\section{GNN scoring siempre 0}

\textbf{Síntoma}: el componente \code{gnn_threat} devuelve 0 para
todo path.

\textbf{Causa}: no hay modelo ONNX cargado. Es el comportamiento
esperado por default.

\textbf{Fix si querés activarlo}:

\begin{itemize}
    \item Descargar o entrenar el modelo (ver documentación del
    core GNN).
    \item Apuntar \code{GRAPHHUNTER_GNN_MODEL_PATH} al \filepath{.onnx}.
    \item Aumentar \code{w5_gnn_threat > 0} en los pesos de
    scoring.
\end{itemize}

\section{Hunt devuelve 0 resultados}

\textbf{Posibles causas}:

\begin{itemize}
    \item Hipótesis demasiado restrictiva (predicados que no
    matchean nada). Relajar un predicado a la vez y correr de
    nuevo.
    \item \code{min_score} demasiado alto. Bajar a 0.0 para ver
    todo.
    \item Ventana temporal (\code{window}) menor que la distancia
    real entre eventos. Aumentar.
    \item Tipos de entidad que no aparecen en el grafo. Verificar
    en \emph{Catálogo}.
\end{itemize}

\section{DLQ creciendo}

\textbf{Síntoma}: \emph{Panel \textrightarrow{} DLQ} acumula
eventos.

\textbf{Causa}: mapping mal configurado, evento con campos
faltantes, o invariante violado.

\textbf{Fix}:

\begin{enumerate}
    \item Abrir DLQ, ver la razón de rechazo.
    \item Si es un error de mapping, corregir el FieldConfig y
    reingestar.
    \item Si es un evento puntual malformado, reingerir tras
    corrección o purgar.
\end{enumerate}

\section{La app está lenta}

\begin{itemize}
    \item Grafo grande (\textgreater~millones de nodos): usar
    modo explorador con \code{Expand grouped} en lugar de expand
    plano.
    \item Hunts sin \code{k} ni \code{window} sobre grafos grandes
    toman minutos — acotar.
    \item Ventana temporal de ingesta muy amplia: considerar
    sesiones separadas por día / semana.
\end{itemize}

\section{Dónde reportar bugs}

En \filepath{https://github.com/base4sec/graphhunter/issues}. Si
abrís un bug, incluir:

\begin{itemize}
    \item Versión del app (\emph{Ayuda \textrightarrow{} Acerca
    de}).
    \item Sistema operativo.
    \item Los logs del app relevantes (\emph{Ver \textrightarrow{}
    Logs del sistema}).
    \item Si aplica, un sample mínimo del archivo que dispara el
    problema — nunca datos reales sensibles.
\end{itemize}
